\section*{Development revision and validation status: 29 October 2026}
This is a working draft for author and specialist review. The revision corrects
hidden-state drift accounting, distinguishes cumulative displacement from peak
mark-to-market loss, and restricts claims to their stated model and observation
contracts. It does not declare every theorem independently verified or every
research objective complete.

\paragraph{Returned execution record.} The author returned a completed full-mode
notebook archive using source commit \texttt{19805d7e}, seed 20261029,
Python 3.13.18 and NumPy 2.1.3. All 10 exported-file hashes and the four analytical
source modules were checked. The run completed E1, E2, E3, E4 and E6; its log
reports 18 targeted tests passing. This is an execution reproduction of the same
seeded code, not independent scientific validation. No real-market CSV was
provided, and the GPU, known-source-feedback and prospective studies were not run.

\paragraph{What the checks establish.} E1 cross-checks the corrected drift identity
against a resolvent and exercises a sufficient invariant-box guard. E2 demonstrates
filter-induced persistence under martingale nulls. E3 compares known alternatives
with equal aggregate risk; it is an illustrative witness, not a uniform rate
theorem. E4 matches baseline capacity on training episodes and reports test
capacity differences; the guard and volatility control both have zero observed
budget breaches in the three bounded-shock synthetic regimes. This does not
establish universal guard superiority or economic welfare.

\paragraph{Coverage and the observation contract.} E6 uses 1,000 replications,
32 reset responses and 512 taps. The nominal 95\% interval assumes independent
Gaussian tap errors and a valid memory envelope. The intervals in the table
measure Monte Carlo uncertainty about coverage, not financial risk.
\begin{center}\small
\begin{tabular}{lrrc}
\toprule
Setting & Coverage & MC 95\% interval & Contract holds\\
\midrule
Independent Gaussian & 94.6\% & [93.0, 95.8] & yes\\
Correlated tap errors & 48.3\% & [45.2, 51.4] & no\\
Heavy-tailed readings & 95.1\% & [93.6, 96.3] & no\\
Biased readings & 25.6\% & [23.0, 28.4] & no\\
Unseen slow memory & 0.0\% & [0.0, 0.4] & no\\
\bottomrule
\end{tabular}
\end{center}
Failures under correlated errors, biased readings and a false memory floor expose
the contract's limits. The heavy-tail outcome in this configuration is not a
general robustness result. All outcomes, including negative ones, are retained.

\paragraph{Remaining obligations.} Independent review of the load-bearing
identification and certification proofs, continuous numerical extrema, economic
primitives, the public-observation bridge and policy counterfactuals remains
necessary. The current two-paper name-map audit finds 54 of 57 formal environments
with existing declaration mappings, three unmapped statements and seven revised
statements requiring semantic review. A label mapping is not a new Lean proof.
The experiment design and full results are retained under
\texttt{studies/economic\_review\_20261029/}.
